\chapter{Traces, Agents and Independence}\label{ch:traces}

A record of an argument must say where the argument came from. The question is what ``where'' can mean. The obvious answer names an author, and the obvious reading of the author is an agent with beliefs and intentions. This chapter argues that the specification can do without that reading, and that doing without it is what lets it handle automated contributors, collaborations and duplicated origin in one framework.

\section{Attribution as an unresolved trace}\label{sec:attribution}

To say that an agent produced a contribution is to say that something upstream of the contribution, unobserved in detail, is being treated as one source. The word ``agent'' compresses a history: prompts, tools, training, edits, reviews, the decision to submit. For a person the compression is habitual and usually harmless. For a system of automated contributors it is not: ten instances of one model, given the same instruction and the same sources, are ten records but very nearly one origin, and any accounting that counts them as ten has counted the compression boundary and not the sources behind it.

The specification therefore records \emph{traces} and not intentions. A trace of an act is what can be inspected about how it came to be: the records it refers to, the sources it used, the method applied, the production mode, the actor identifier, and the sequence position. Trace completeness has six components, and an act with any component missing is trace-null (Definition in Chapter~\ref{ch:ledger}). A trace-null act is not an act by an unknown agent; it is an act whose origin is unresolved, and the ledger says so. Such an act is quarantined and inert: it is stored, addressable and citable as a record, but takes no part in evaluation, and adds nothing to any count.

\begin{designrule}[Unresolved origin is a state, not a default]\label{rule:tracenull}
An act whose trace is incomplete is recorded as such and is never assigned a default origin. Assigning one (a default human author, say) would convert an absence into an assertion, and later counts of independent sources would depend on it.
\end{designrule}

\section{Independence is a property of origins}\label{sec:originindep}

Evidence counts as independent support only when it has independent origins, and origin is determined by shared methodological dependence, not by the number of records. Chapter~\ref{ch:evidence} defines independence families as classes of a signed identification and proves that duplicates do not raise the independent count (Proposition~\ref{prop:indep}). The consequences for automated and imported material are direct.
\begin{itemize}
\item A batch of records generated by one procedure from one source is one family, however many records it contains.
\item Repeated runs of one model with one instruction over one corpus are one family unless the runs differ in a way that a declared equivalence claim, itself open to objection, recognizes as producing independent evidence.
\item Agreement among records of one family is displayed, since it is information, but never counted as replication.
\end{itemize}
The rule is deliberately conservative. A false merge of two independent sources understates support and is repaired by a distinction claim; a false separation of two dependent sources overstates support and is repaired by an equivalence claim. Both repairs are ordinary acts in the ledger, contestable like any other, and both are visible in the status vector through the coordinate that counts families.

\section{Inspectable work and performed agency}\label{sec:inspectable}

A public record can reward the appearance of purposeful activity, such as many posts, many replies and many actors, or it can reward acts whose reasons can be checked. The specification takes the second course. Credit, standing and influence attach to argumentative acts, and an act counts only through the record of what it did to the argument graph. Goal-directedness is attributed by observers and is not stored. This has one practical consequence that matters for automated contributors: a system that produces a great deal of fluent activity gains nothing from the volume, since volume neither adds independent families (Section~\ref{sec:originindep}) nor carries any defeat power (Proposition~\ref{prop:agreement}), and gains nothing from apparent agency, since agency is not an input to evaluation (Proposition~\ref{prop:separation}).

What an automated contributor can do is what any contributor can do: add evidence with a checked anchor, supply a warrant with backing, raise an objection whose active region is nonempty, or restrict an overbroad claim. Each of these is a legible act with an effect on the graph. The specification does not decide whether such an act is wise. It decides that the act is recorded with its origin mode and evaluated by the same rules as every other.
